% ============================================================
%  appendix/grammar.tex
%  Appendix C — BNF Grammar for Grade Notation and Op1–Op10
%  STATUS: Write W1
%  Note: reflect Op10 domain restriction (OPEN) per §3.
% ============================================================

% Note: this file is \input{} from main.tex inside \appendix.

\TODO{Write complete BNF grammar. Template below is a starting point.}

\medskip
The following BNF grammar defines valid grade expressions and operation
applications in the E1--E7 extension stack.

\begin{verbatim}
<grade>       ::= "(" <nat> "," <nat> ")"
<nat>         ::= "0" | "1" | "2" | ...

<element>     ::= "VACUUM"
               | "F" "(" <nat> "," <nat> ")"

<non-vacuum>  ::= "F" "(" <nat> "," <nat> ")"   -- (n,m) ≠ (0,0)

<op-name>     ::= "Op1" | "Op2" | "Op3" | "Op4" | "Op5"
               | "Op6" | "Op7" | "Op8" | "Op9" | "Op10"

<application> ::= <op-name> "(" <arg-list> ")"

<arg-list>    ::= <element>
               | <element> "," <arg-list>

-- Op-specific domain constraints:
-- Op2 : <element> → <non-vacuum>     (α grade-shift; codomain ≠ VACUUM for n≥1)
-- Op3 : <element> → <non-vacuum>     (β grade-shift; codomain ≠ VACUUM)
-- Op5 : <non-vacuum> → <non-vacuum>  (inversion; VACUUM has no inverse)
-- Op9 : Ω_k → op-subset             (threshold values VALIDATED;
--                                     selection mechanism CONJECTURE)
-- Op10: <triple> → {0,1}            on individual triples  [DEFINED]
--       Op10(Op(x))                  iterated form requires
--                                     domain closure of Op9  [OPEN]
\end{verbatim}

\begin{remark}
Op10 domain restriction is labelled \epistemic{OPEN}: the iterated form
$\mathrm{Op10}(\mathrm{Op}(x))$ requires that Op9 closes the domain under the
relevant operation. This condition has not been analytically verified.
\end{remark}

\TODO{Expand grammar with explicit production rules for each operation.
Cross-reference to Axioms~1--4 and Table~\ref{tab:ops}.}
